In chapter \ref{ch:proc-pascal} we learned how to declare and call procedures without parameters,
which let us repeat the same sequence of commands of the
\robot\ several times at different places in the program. In chapter \ref{ch:appel} we
used commands of the \robot\ which are calls of procedures with parameters,
the use of one or more parameters al\-low\-ing the procedure to perform sequences
of commands that are similar but differ according to the value of the parameters. We will now
see how to declare procedures with parameters in \python. Let us start
with a simple example:

\bigskip

 {\raggedright 
 {}\hfill\vbox to 0pt{\vspace*{\baselineskip}\vspace*{3cm}
 \includevspacePostScript[-\baselineskip]{CarrePointille}{0.3}{0.3}}\hspace*{30mm}%
 \label{prog:CarrePointille}
 \input python/DottedSquare.tex
 }
The procedure {\idt dottedLine} draws a dotted line made up of a succession of
{\idt n} black pixels separated by $(${\idt n}$ - 1)$ white pixels, the num\-ber {\idt n} of dots being
passed as a parameter of this procedure. In the declaration `{\mcl def}\ {\idt dottedLine(n):}',
the parameter {\idt n} is written between the parentheses that follow the name of the procedure,
parentheses which remain empty {\idt ()} for procedures without parameters. There is no need to state that {\idt n} is
an integer: \python\ does not ask for it, and it is only at execution time,
when computing the expression {\idt n - 1}, that it uses the value of {\idt n}.
The procedure {\idt square} draws a dotted
square, the number {\idt p} of dots per side being passed as a parameter. The call {\idt square(30)}
says that the procedure {\idt square} must be executed with a parameter {\idt p} equal to 30. Consequently the
procedure {\idt dottedLine} is called four times with a parameter {\idt n} equal to {\idt p}, that is
30. We say that {\idt p} is the {\em formal parameter}\index{Formal parameter}  (whose value is not yet known
at the time the procedure is declared, and which is designated by an identifier) and that 
30 is the {\em actual parameter}\index{Actual parameter} (that is,
the value of the formal parameter for this call of the procedure).

In general, declarations and calls of procedures with pa\-ram\-e\-ters have the following
form:

\index{Formal parameter}\index{Procedure declaration}\index{Identifier}
\begin{center}
\begin{tabular}{@{}|p{12.5cm}|@{}} \hline
{\em Declaration of a procedure with parameters:}\\ \hline \hline
\begin{minipage}[t]{12.5cm}
{\medskip
\mbox{\blc\ \ \ \cmt \#\ \small\rm Specification: what the procedure does}\par 
\mbox{\blc\ \ \ \mcl def\blc\hspace{0.3em}\blanc\hspace{0.3em}\idt procedureName(%
formalParam1\footnote{\ A formal parameter is an identifier.},\ formalParam2,\ \dots):}\par 
\mbox{\blc\ \ \ \ \ \ \ \small\rm \ldots\ Instructions of the body of the procedure%
\footnote{\ The expressions appearing in these instructions may use the formal parameters
{\idt formalParam1}, \ldots\ of the procedure.}
 (indented) \ldots}\par 
\medskip}
\end{minipage}\\ \hline 
\end{tabular}
\end{center}

\index{Actual parameter}\index{Procedure call}
\begin{center}
\begin{tabular}{@{}|p{12.5cm}|@{}} \hline
{\em Call of a procedure with parameters:}\\ \hline \hline
\begin{minipage}[t]{12.5cm}
{\medskip
\mbox{\blc\ \ \ \idt procedureName(actualParam1%
\footnote{\ An actual parameter is an integer expression (or a colour, etc.).},\ actualParam2,\ \dots%
\footnote{\ The actual parameters are given in the order in which the corresponding formal
pa\-ram\-e\-ters appear.})}\par 
\medskip}
\end{minipage}\\ \hline 
\end{tabular}
\end{center}

Note that the parameters, formal ones in the procedure declaration as well as actual ones
in the procedure call, are placed between parentheses and separated by
commas. For example, to draw a rectangle of height 3 and width 4 with the 
procedure {\idt rectangle} declared by `{\mcl def}\ {\idt rectangle(height,\ width):}' in the program
{\idt MacintoshSE.py} of the folder 
{\idt book/\allowbreak 13\_declaring\_procedures}, we write {\idt rectangle(3,\ 4)}.
If we forget the parentheses and write {\idt rectangle\ 3,\ 4}, \python\ rejects the program
before running it with the message `{\tt SyntaxError: invalid syntax}'.
An actual parameter can be an expression,
as in {\idt rightHalfCircle(2\ *\ gridSize)} (page \pageref{prog:HeliceEnCouleur})
or {\idt fdn(-3)}: it is computed before the procedure is called. Be careful to place
the closing parenthesis correctly: {\idt rightHalfCircle(2)\ *\ gridSize} would be
understood as the multiplication by {\idt gridSize} of the result of the call {\idt rightHalfCircle(2)}, 
which makes no sense. After drawing the half circle, \python\ stops with the message:
`{\tt TypeError: unsupported operand type(s) for *: 'NoneType' and 'int'}'
(operand types not allowed for {\idt *}: {\idt NoneType} and integer). {\idt NoneType} is the type
of the value {\mcl None} (`nothing') returned by a procedure call.

The number of actual parameters must be equal to the number of formal parameters.
If we write {\idt rectangle(3)}, \python\ stops when the \robot\ reaches this call,
with the message
`{\tt TypeError: rectangle() missing 1 required positional argument: 'width'}'
and if we write
{\idt rectangle(3,\ 4,\ 5)}, with the message
`{\tt TypeError: rectangle() takes 2 positional arguments but 3 were given}'.


A procedure with 
formal parameters can be called several times in a program with 
actual pa\-ram\-e\-ters
that can be different each time.

The formal parameters of a procedure are names that are {\em local} to the procedure: they can only be used in the
body of the procedure (elsewhere, \python\ would answer for example
`{\tt NameError: name 'width' is not defined}'). Consequently different procedures can have
formal parameters with the same name although they are different parameters. 

These various 
points are illustrated by the example below.

\noindent\begin{minipage}{\textwidth}
\begin{exemple}[Segmented spiral]{\rm
~The program below draws a spiral made up of pairs of straight line segments
of successive lengths 1, 2, 3, \ldots, 19:

\nopagebreak
 {\raggedright 
 \mbox{}\hspace*{40mm}\vbox to 0pt{\vspace*{\baselineskip}\vspace*{4cm}
\includevspacePostScript[-\baselineskip]{SpiraleSegmentee}{0.7}{0.7}}
 \input python/SegmentedSpiral.tex
 }
\noindent $\Box$ }
\end{exemple}
\end{minipage}


\needspace{6\baselineskip}\noindent \rule{137mm}{0.1mm}\par\nopagebreak %---------------------------------------------
\begin{exercices}{Write \python\ programs containing procedures with
parameters to produce the following drawings:

\tiret the cross with three transverse arms, the coloured propeller, the mi\-cro\-com\-put\-er,
the colour logo, the caterpillar (where the number
of rings of the caterpillar's body is passed as a parameter of the drawing procedure), the canoe 
(where the number of rowers is a parameter) shown on page \pageref{fig:crx3b};

\tiret the squares inscribed one inside another
or the painted version of this drawing shown in figure \ref{fig:carinspt} (page \pageref{fig:carinspt});

\tiret the arabesque inspired by the dome of the mausoleum of Sultan
         Qayt-\-Bay shown in figure \ref{fig:arab2} (page \pageref{fig:arab2});

\tiret the barge of figure \ref{fig:transports} (page \pageref{fig:transports}) where the number
         of portholes is given as a parameter;

\tiret the Golomb tiling of figure \ref{fig:golomb} (page \pageref{fig:golomb}).

\noindent $\Box$}\end{exercices}
\twocolumn
\begin{figure}[h]
%\hspace{-3.363cm} 
\includevspacePostScript[4.9cm]{Arabesque2}{0.95}{0.95}
 \caption{Arabesque}\label{fig:arab2}
\end{figure}
\vfill
\begin{figure}[h]
%\hspace{-3.363cm} 
\includevspacePostScript[0.0cm]{HeliceEnCouleur}{0.95}{0.95}
 \caption{Coloured propeller}
 %\label{fig:helice}
\end{figure}
\vfill
\begin{figure}[h]
%\hspace{-3.363cm} 
\includevspacePostScript[0.0cm]{MacintoshSE}{0.95}{0.95}
 \caption{Old microcomputer}
 %\label{fig:MacSE}
\end{figure}
\mbox{}
\newpage
\begin{figure}[!h]
%\hspace{-3.363cm} 
\includevspacePostScript[5.3cm]{LogoEnCouleur}{0.95}{0.95}%
\caption{Colour logo}%
%\label{fig:Logoc}
\end{figure}
\begin{figure}[!h]
%\hspace{-3.363cm} 
\includevspacePostScript[1.0cm]{Chenille}{0.95}{0.95}%
\caption{Caterpillar}%
%\label{fig:chenille}
\end{figure}
\begin{figure}[!h]
%\hspace{-3.363cm} 
\includevspacePostScript[1.0cm]{Pirogue_1}{0.95}{0.95}%
\caption{Canoe}%
%\label{fig:pirog1}
\end{figure}
\vfil
\begin{figure}[!h]
\includevspacePostScript[0.0cm]{Croix3branches}{0.95}{0.95}%
\caption{Cross with three arms}%
\label{fig:crx3b}%
\end{figure}
\onecolumn
\begin{figure}[h]
\includevspacePostScript[5.7cm]{CarresInscritsPeints}{0.95}{0.95}
 \caption{Horemis's pattern of inscribed squares}\label{fig:carinspt}
\end{figure}
\begin{corrige}{
We give one program, the others being in the folder 
{\idt book/\allowbreak 13\_declaring\_procedures}:
\label{prog:HeliceEnCouleur}
\input python/ColourPropeller.tex
\noindent $\Box$}\end{corrige}

\vfill

%%%%%%%%%%%%%%%%%%%%%%%%%%%%%%%%%%%%%%%%%%%%%%%%%%%%%%%%%%%%%%%%%%%%%%%%%%%
\noindent \framebox[\textwidth]{
\begin{minipage}{\myparboxwidth}
\setlength{\parindent}{2em}
\small \sf
\vspace{\abovedisplayskip}
\section*{Addenda}

In \python, the definitions of identifiers are distinguished from their uses. \index{Identifier}
The   {\em defining occurrences} appear in a constant
or procedure declaration ({\mcl def}) as the name of a constant, of a procedure or of a formal parameter.
The  {\em using occurrences} of these identifiers are the names of constants or of
formal parameters used in an integer expression, or the names of procedures used in calls.

 The same identifier may have several defining occurrences in a program. Hence the
 problem arises of matching the using occurrences with the defining occurrences of the same
 identifier. For 
example the identifier can be used to name a constant of the program and formal parameters
in procedures. In \python, these procedures can be nested one inside
 another (by declaring with {\mcl def} a procedure {\idt p2} inside the body of a procedure {\idt p1},
 with one more level of indentation), as in the example below:
 
\smallskip

{
\parindent \parindentprog
\medskip

\vspace{\neglskip}\mbox{\blc {\mcl from}{\idt \ LibRobot\ }{\mcl import}{\idt \ *}}\par
\vspace{0.3\baselineskip}
\vspace{\neglskip}\mbox{\blc {\idt n\ =\ 17}}\par
\vspace{0.3\baselineskip}
\vspace{\neglskip}\mbox{\blc {\mcl def}{\idt \ p1(n):}}\par
\vspace{\neglskip}\mbox{\blc \blk\blk{\mcl def}{\idt \ p2(n):}}\par
\vspace{\neglskip}\mbox{\blc \blk\blk\blk\blk{\idt ...}}\par
\vspace{\neglskip}\mbox{\blc \blk\blk\blk\blk{\idt ...\ n\ ...\ \ }{\cmt \#1}}\par
\vspace{\neglskip}\mbox{\blc \blk\blk\blk\blk{\idt ...}}\par
\vspace{\neglskip}\mbox{\blc \blk\blk{\idt ...}}\par
\vspace{\neglskip}\mbox{\blc \blk\blk{\idt p2(n)\ \ }{\cmt \#2}}\par
\vspace{\neglskip}\mbox{\blc \blk\blk{\idt ...}}\par
\vspace{0.3\baselineskip}
\vspace{\neglskip}\mbox{\blc {\idt ...}}\par
\vspace{\neglskip}\mbox{\blc {\idt p1(n)\ \ }{\cmt \#3}}\par
\vspace{\neglskip}\mbox{\blc {\idt ...}}\par

\medskip }

We say that the program has a {\em block structure}. \index{Block structure}These blocks are formed by the procedures nested 
 one inside another, the program being considered as the outermost block. 
 
 The {\em scope 
 rule}\index{Scope} of identifiers in \python\ says that a using occurrence of an identifier
 refers to the defining occurrence of the innermost block in which this using occurrence appears and
 which contains a defining occurrence of the identifier. 

\vspace{\abovedisplayskip}
\end{minipage}
}

%%%%%%%%%%%%%%%%%%%%%%%%%%%%%%%%%%%%%%%%%%%%%%%%%%%%%%%%%%%%%%%%%%%%%%%%%%%% %%%%%%%%%%%%%%%%%%%%%%%%%%%%%%%%%%%%%%%%%%%%%%%%%%%%%%%%%%%%%%%%%%%%%%%%%%%
%%%%%%%%%%%%%%%%%%%%%%%%%%%%%%%%%%%%%%%%%%%%%%%%%%%%%%%%%%%%%%%%%%%%%%%%%%%
\noindent \framebox[\textwidth]{
\begin{minipage}{\myparboxwidth}
\setlength{\parindent}{2em}
\small \sf
\vspace{\abovedisplayskip}
 
 
In the example above,
\begin{list}{$\scriptscriptstyle\bullet$}{
}
 \item the using occurrence of {\idt n} on the line marked {\cmt \#1} refers to the formal parameter of the procedure {\idt p2};
\item the using occurrence of {\idt n} on the line marked {\cmt \#2} refers to the formal parameter of the procedure {\idt p1};
 \item the using occurrence of {\idt n} on the line marked {\cmt \#3} refers to the constant 17.
\end{list}

Block structure, which is found in \python, was introduced by {\sc algol 60} \cite{algol60}. This language was created by a group
of European and American experts gathered in the working group WG2.1 of IFIP (International Federation for 
Information Processing), which brings together all the national computing societies. It is the ancestor of a family
of languages which includes \pascal, {\sc Ada} and {\sc Modula}.

\vspace{\abovedisplayskip}
\end{minipage}
}
%%%%%%%%%%%%%%%%%%%%%%%%%%%%%%%%%%%%%%%%%%%%%%%%%%%%%%%%%%%%%%%%%%%%%%%%%%%%
